\section{Introduction}\label{sect_introduction1}%
%

Lattice calculations in QCD have demonstrated the ability to complement, and in certain cases with the exceeding precision, significant amount of experimental measurements.  Now, the lattice evaluation of parton distribution functions (PDFs) are coming on the agenda. New techniques are being explored aiming at the access of PDFs directly in the momentum fraction space, in addition to the standard approach that allows one to calculate first Mellin moments of PDFs. The existing actual proposals~\cite{Aglietti:1998ur,Abada:2001if,Detmold:2005gg,Braun:2007wv,Ji:2013dva,Ma:2017pxb} differ in details but have a common general scheme: PDFs are extracted from the lattice calculations of suitable Euclidean correlation functions using QCD collinear factorization in the continuum theory. 

A particularly popular suggestion~\cite{Ji:2013dva} that has triggered  a lot of recent activity \cite{Lin:2014zya,Alexandrou:2015rja,Chen:2016utp,Alexandrou:2016jqi,Lin:2017ani,Zhang:2017bzy,Chen:2017gck,Chen:2017mzz,Alexandrou:2017huk, Alexandrou:2018pbm,Chen:2018xof,Chen:2018fwa,Alexandrou:2018eet,
Lin:2018qky}, introduces a concept of a parton \emph{quasi-distribution} (qPDF) defined as a Fourier transform of the nonlocal quark-antiquark operator connected by the Wilson line. QPDF is then matched to PDF either directly in the $\MSbar$ scheme or using the large-momentum factorization scheme at the intermediate step (``Large Momentum Effective Theory'' (LaMET)~\cite{Ji:2014gla,Izubuchi:2018srq}). The latter technique is useful to emphasize that, after the Fourier transform, the hadron momentum $p$ remains to be the only dimensional parameter. So, the relevant scale of the QCD coupling is related to $p$ up to a dimensionless constant and also the higher-twist corrections are generically suppressed by powers of the momentum $p$. A closely related quantity, a \emph{pseudo-distribution} (pPDF) was introduced in~\cite{Radyushkin:2017cyf,Orginos:2017kos,Karpie:2017bzm}.
 
Lattice calculations of PDFs using these new methods are presently moving from exploratory stage towards precision calculations, therefore questions like whether the higher-twist (power suppressed) corrections are well under control have to be addressed. One possibility to investigate the impact of higher-twist corrections is to extract PDFs from the global analysis of many Euclidean correlation functions introducing 
{such} corrections as free parameters. This approach is
showcased in \cite{Bali:2018spj} using the position-space strategy~\cite{Braun:2007wv,Bali:2017gfr} 
and the pion light-cone distribution amplitude (LCDA) as an example.
%
In this case, the analytic structure of higher-twist effects is well understood, so they could have been modelled by just one free parameter.  
The general situation  may be more complicated. Having in mind the multitude of 
observables that can potentially be employed in lattice calculations, see e.g.~\cite{Ma:2017pxb}, 
it is desirable to have a general method to estimate the corresponding higher-twist corrections and their expected Bjorken-$x$ dependence.   
The purpose of this work is to point out that the problem of power-suppressed contributions for such observables 
can be addressed using the concept of renormalons \cite{Beneke:1998ui,Beneke:2000kc}. In what follows we apply this technique 
to qPDFs and pPDFs.

The renormalon approach to the investigation of power corrections is founded on the fact that operators of different twist mix with each other under renormalization, due to the violation of QCD scale invariance through the running of the coupling constant. In cutoff schemes, this mixing is explicit, whereas in dimensional regularization, it manifests itself in factorial divergence of the perturbative series. Independence of a physical observable on the factorization scale implies intricate cancellations between different twists --- the cancellation of renormalon ambiguities. In turn, the existence of these ambiguities in the leading-twist expressions can be used to estimate the size of power-suppressed corrections. Conceptually, it is similar to the estimation of the accuracy of fixed-order perturbative results by the logarithmic scale dependence. The renormalon approach was used before for the study of Bjorken-$x$ dependence of higher-twist corrections in deep-inelastic scattering \cite{Beneke:1995pq,Dokshitzer:1995qm,Dasgupta:1996hh}, fragmentation functions~\cite{Dasgupta:1996ki,Beneke:1997sr}, pion LCDA~\cite{Braun:2004bu} and transverse-momentum-dependent parton distributions~\cite{Scimemi:2016ffw}.

In order to explain how the concept of renormalons can be used to get insight in the structure of power corrections, let us
consider the usual expression for the quasi-distribution~\cite{Ji:2015qla,Izubuchi:2018srq}, 
\begin{align}
\label{qPDF1}
 \mathcal{Q}(x,p) &= \int_{-1}^1 \!\frac{dy}{|y|} C_{\Qd}(\tfrac{x}{y},xp,\mu_F) q(y,\mu_F) + \frac{1}{p^2}  \mathcal{Q}_4(x,p) + \ldots
\end{align}
where $q(y,\mu_F)$ is the quark PDF, $p$ is the hadron momentum, $x$ refers to the momentum fraction. For brevity we do not show the dependence on the renormalization scale. The factorization scale $\mu_F$ has to be taken of the order of $|x| p$ to avoid large logarithms.  The coefficient function $C(x,p,\mu_F) = \delta(1-x) + \mathcal{O}(\alpha_s)$ is given by the perturbative expansion. The correction of ${\cal O}(\alpha_s)$ was first computed for the flavor-nonsinglet case in \cite{Xiong:2013bka}, see also~\cite{Stewart:2017tvs}. 

To understand the role of renormalons, it is necessary to examine carefully the separation made in~\eqref{qPDF1} between the leading term and the higher twist addenda. Let us assume for a moment that the factorization is done using a hard cutoff $\Lambda_{\rm QCD} \ll \mu_F  \ll p$, i.e. the contributions with loop momenta $|k| > \mu_F$ are included  in the coefficient function, whereas the contributions with $|k| < \mu_F$ are included in PDF. In this scheme, 
the coefficient function has the following expansion at $p\to\infty$
\begin{align}
\label{C-cutoff}
 C_{\Qd}(x,p,\mu_F) &= \delta(1-x) + c_1 \alpha_s + c_2\alpha_s^2 + \ldots 
\notag\\ &\quad{}
- \frac{\mu_F^2}{p^2} D_{\Qd}(x) +\ldots\,,
\end{align}    
where $c_k= c_k(x,\ln p^2/\mu_F^2)$ are the perturbative coefficients depending logarithmically on the scales and the $D_{\Qd}$-term represents the leading power correction. Since the left-hand-side of Eq.~\eqref{qPDF1} does not depend on $\mu_F$, any such dependence should cancel on the right-hand-side. In particular, the logarithmic dependence on the scale in $c_k(x,\ln p^2/\mu_F^2)$ is canceled by the scale-dependence of PDF $q(x,\mu_F)$. The cancellation of the power dependence, on the other hand, must involve the twist-four contribution $\mathcal{Q}_4(x,p)$. Thus, in this factorization scheme one expects that
\begin{align}
\label{Q4-cutoff}
 \mathcal{Q}_4(x,p,\mu_F) = \mu_F^2 \int_{-1}^1 \!\frac{dy}{|y|} D_{\Qd}(\tfrac{x}{y}) q(y,\mu_F) + \widetilde{\mathcal{Q}}_4(x,p,\mu_F)\,,  
\end{align} 
where $\widetilde{\mathcal{Q}}_4$ depends on $\mu_F$ at most logarithmically. 
Appearance of the term $\sim \mu_F^2$ can be traced to quadratic ultraviolet divergence (in addition to the logarithmic ultraviolet divergence) 
of the twist-four operators that are responsible for the power correction, in the cutoff scheme. 
One can prove that the cutoff dependence $\sim \mu_F^2$ of the higher-twist operators is indeed that of Eq.~\eqref{C-cutoff}.    

In practice, perturbative calculations are usually done using dimensional regularization. In this case, power-like terms as in~\eqref{C-cutoff} do not appear. The price to pay is that the coefficients $c_k$ computed in a $\MSbar$ scheme grow factorially with the order $k$. The factorial growth implies that the sum of the perturbative series is only defined to a power accuracy and this ambiguity (renormalon ambiguity) must be compensated by adding a non-perturbative higher--twist correction. The detailed analysis shows~\cite{Beneke:1998ui,Beneke:2000kc} that the divergent large-$k$ behavior (the renomalon) of the coefficients is in the one-to-one correspondence with the sensitivity to extreme (small or large) loop momenta. In particular, infrared renormalons in  the leading-twist coefficient function are compensated by ultraviolet renormalons in the matrix elements of twist-four operators. In this way the same picture as in the cutoff scheme re-appears in dimension regularization.

Returning to \eqref{Q4-cutoff}, we observe that the quadratic term in $\mu_F$ is spurious since its sole purpose is to cancel the similar contribution to the coefficient function. Therefore, it does not contribute to any physical observable. The idea of the renormalon model of the power corrections \cite{Beneke:1995pq,Dokshitzer:1995qm,Dasgupta:1996hh,Dasgupta:1996ki,Beneke:1997sr,Braun:2004bu} is that, with a replacement of $\mu_F$ by a suitable non-perturbative scale, this contribution reflects the order and the functional form of actual power-suppressed contribution. Assuming this ``ultraviolet dominance''~\cite{Braun:1995tp,Beneke:1998ui,Beneke:2000kc} one obtains the following  model:
\begin{equation}
 \mathcal{Q}_4(x,p,\mu_F) = \kappa 
\Lambda_{\rm QCD}^2 \int_{-1}^1 \!\frac{dy}{|y|} D_{\Qd}(\tfrac{x}{y})\, q(y,\mu_F)\,,
\label{UV_Lambda_int}
\end{equation}
with the dimensionless coefficient $\kappa = \mathcal{O}(1)$ which cannot be fixed within theory and remains a free parameter.

In this work we calculate the function $D_{\Qd}(x)$ for different versions of the qPDFs (pPDFs) in the so-called bubble-chain approximation~\cite{Beneke:1998ui}, which is our main result. This calculation reveals that the power corrections to quasi- and pseudo-distributions have the
following generic structure
\begin{align}
 \cQ (x,p) &=  q(x) \Big\{1 + \mathcal{O}\Big( \frac{\Lambda^2}{p^2x^2(1-x)}\Big) \Big\},  
\notag\\
\mathcal{P}(x,z) &= q(x)\Big\{1 + \mathcal{O}\big( z^2\!\Lambda^2{(1-x)}\big)\Big\},
\end{align}
respectively, and can be affected significantly by normalization. In particular, the normalization of the involved matrix elements to their value at zero momentum considerably reduces the power correction for the qPDFs at smaller-$x$ at the cost of a strong enhancement at larger values. We emphasize that the leading power correction to qPDF at $x\to 0$  is enhanced by two powers of the Bjorken $x$ variable. For pPDFs the power corrections are suppresed at $x\to1$, which, unfortunately,  does not hold after the normalization procedure of Ref.~\cite{Orginos:2017kos} (but can be upheld with a different choice). Additionally, as a byproduct of bubble-chain calculation, we have obtained the large-$n_f$ part, $n^k_f \alpha_s^{k+1}$, of the
leading twist coefficient function to all orders in perturbation theory.   

    

The presentation is organized as follows. In Sect.~2 we formulate
our program in more precise terms. We define qPDFs and  pPDFs as 
particular Fourier transforms of the position space (Ioffe-time) distributions, 
discuss briefly the light-ray operator product expansion (OPE) and the target mass corrections, and  
introduce the relevant techniques (Borel transform) and the systematic approximation
(large-$n_f$ expansion) that will be used throughout the rest of
the work. In Sect. 3 we present our result for the Borel transform of the leading-twist
coefficient function and discuss the structure of its singularities. 
The leading power corrections to various versions of the quasi-distributions are obtained in Sect.~4. 
We collect there the relevant analytic expressions and also present the results of a numerical 
study using realistic parametrizations for the valence quark PDFs. 
The final Sect.~5 is reserved for the summary and conclusions.   


%
